\documentclass[12pt, a4paper]{article}
\usepackage[ ]{indentfirst}
\usepackage[brazilian]{babel}
\title{Revisão P1\\Engenharia de Software II}
\author{Erickson}
\date{\today}

\begin{document}
	\maketitle
	\newpage
	
	\section{Diagrama de Casos de Uso}
		Cada requisito é um caso de uso. São representados no diagrama por uma elipse.
		\subsection{Atores}
			Ator é alguém ou algo que interage com o sistema.
	\section{Diagrama de Classes}
		Tem como objetivo trazer como os daods vão ser armazenados dentro do programa. O diagrama de classes só pode existir depois da criação do diagrama de casos de uso.
		\subsection{Classe}
			Exemplo: Carro

			Atributos: marca, modelo,cor

		\subsection{Objeto}
			Exemplo da instância do objeto: fiat, Uno, Branco

		\subsection{Método}
			Dirigir()
		
\end{document}
